\documentclass[12pt,a4paper]{article}

% --- 1. Configuración de Idioma y Formato ---
\usepackage[utf8]{inputenc}
\usepackage[spanish]{babel}
\usepackage[margin=2.5cm]{geometry}
\usepackage{xcolor}
\usepackage{graphicx}
\usepackage[T1]{fontenc}
\usepackage{hyperref}
\usepackage{titlesec}
\usepackage{booktabs} % Para tablas elegantes
\usepackage{fancyhdr} % Para encabezados y pies de página
\usepackage{lmodern}
\usepackage{microtype}
\usepackage{caption}
\captionsetup[figure]{labelfont=bf, textfont=normalfont, labelsep=colon}

% --- Configuración de Bibliografía ---
\usepackage[backend=biber, style=ieee]{biblatex} % Estilo APA (o usa 'numeric')
\addbibresource{memoria.bib} % Tu archivo .bib

% --- 2. Configuración para Código Fuente (YAML/Bash) ---
\usepackage{listings}
\definecolor{codegreen}{rgb}{0,0.6,0}
\definecolor{codegray}{rgb}{0.5,0.5,0.5}
\definecolor{backcolour}{rgb}{0.96,0.96,0.94}

\setlength{\parskip}{0.7em}

\addtolength{\skip\footins}{20pt}



\lstdefinestyle{techstyle}{
    backgroundcolor=\color{backcolour},   
    commentstyle=\color{codegreen},
    keywordstyle=\color{blue},
    numberstyle=\tiny\color{codegray},
    stringstyle=\color{orange},
    basicstyle=\ttfamily\small,
    breaklines=true,                 
    captionpos=b,                    
    numbers=left,                    
    showstringspaces=false,
    tabsize=2,
    literate={á}{{\'a}}1 {é}{{\'e}}1 {í}{{\'i}}1 {ó}{{\'o}}1 {ú}{{\'u}}1 {ñ}{{\~n}}1
}
\lstset{style=techstyle}

% Define YAML language for listings
\lstdefinelanguage{yaml}{
    backgroundcolor=\color{gris-claro},
    frame=single,                % Caja cerrada sencilla
    rulecolor=\color{gris-borde}, % Color del borde
    basicstyle=\ttfamily\small,  % Fuente monoespaciada
    breaklines=true,             % Cortar líneas largas
    xleftmargin=10pt,            % Margen izquierdo
    xrightmargin=10pt,           % Margen derecho
    showstringspaces=false,
    keywordstyle=\color{blue},   % Comandos en azul
    commentstyle=\color{gray},   % Comentarios en gris
    captionpos=t,                % Título arriba
    literate={á}{{\'a}}1 {é}{{\'e}}1 {í}{{\'i}}1 {ó}{{\'o}}1 {ú}{{\'u}}1 {ñ}{{\~n}}1 {Á}{{\'A}}1 {É}{{\'E}}1 {Í}{{\'I}}1 {Ó}{{\'O}}1 {Ú}{{\'U}}1
}

\definecolor{gris-claro}{rgb}{0.97,0.97,0.97}
\definecolor{gris-borde}{rgb}{0.8,0.8,0.8}

\lstdefinestyle{cajabasica}{
    backgroundcolor=\color{gris-claro},
    frame=single,                % Caja cerrada sencilla
    rulecolor=\color{gris-borde}, % Color del borde
    basicstyle=\ttfamily\small,  % Fuente monoespaciada
    breaklines=true,             % Cortar líneas largas
    xleftmargin=10pt,            % Margen izquierdo
    xrightmargin=10pt,           % Margen derecho
    showstringspaces=false,
    keywordstyle=\color{blue},   % Comandos en azul
    commentstyle=\color{gray},   % Comentarios en gris
    captionpos=t,                % Título arriba
    literate={á}{{\'a}}1 {é}{{\'e}}1 {í}{{\'i}}1 {ó}{{\'o}}1 {ú}{{\'u}}1 {ñ}{{\~n}}1 {Á}{{\'A}}1 {É}{{\'E}}1 {Í}{{\'I}}1 {Ó}{{\'O}}1 {Ú}{{\'U}}1
}

% --- 3. Encabezado y Pie de Página ---
\pagestyle{fancy}
\fancyhf{}
\lhead{Entregable 1 -- Monitor}
\rhead{\thepage}
\cfoot{Julen García Muñoz}
% --- 4. Información de la Portada ---
\title{
    \vspace{2cm}
    \textbf{Entregable 1 -- Monitor}
}
\author{Julen García Muñoz}
\date{\today}

% --- 5. Inicio del Documento ---
\begin{document}

\maketitle
\vfill
\begin{center}
    \textbf{Sistemas empotrados y ubicuos -- SEU} \\
    Entregable 1 -- Monitor
\end{center}
\newpage

\tableofcontents
\newpage

\section{Módulo Monitor}

\subsection{Descripción General}
El módulo \texttt{monitor.c} implementa un sistema de monitoreo de sensores para una placa STM32F411 con interfaz de LEDs, buzzer y botones. El sistema permite la lectura de dos sensores (LDR para luminosidad y NTC para temperatura), ajuste de niveles de alarma mediante potenciómetro, visualización en LEDs y activación de alarma sonora.

\subsection{Constantes}

\subsubsection{Constantes NTC}
\begin{lstlisting}[language=C, caption=Parámetros del termistor NTC]
#define R25    10000.0f  // Resistencia a 25 deg C en ohmios
#define T25    298.15f   // Temperatura de referencia en Kelvin (25 deg C)
#define BETA   3900.0f   // Coeficiente Beta del termistor
\end{lstlisting}

\subsubsection{Canales ADC}
\begin{lstlisting}[language=C, caption=Asignación de canales ADC]
#define CH_LDR ADC_CHANNEL_0   // Canal para sensor LDR
#define CH_NTC ADC_CHANNEL_1   // Canal para sensor NTC
#define CH_POT ADC_CHANNEL_4   // Canal para potenciómetro
\end{lstlisting}

\subsection{Tipos de Datos}

\subsubsection{Estructura sensor\_t}
\begin{lstlisting}[language=C, caption=Estructura para almacenar datos de sensores]
typedef struct {
    float valor;                  // Valor actual del sensor
    float minimo;                 // Valor mínimo de rango
    float maximo;                 // Valor máximo de rango
    float nivel_alarma;           // Umbral de activación de alarma
    int activated;                // Flag de activación
    uint32_t time_activation;     // Timestamp de activación
    int value_flashing;           // Estado del parpadeo
    uint32_t flashing_last_time;  // Timestamp del último cambio de parpadeo
} sensor_t;
\end{lstlisting}

\subsubsection{Enumeración alarm\_state\_t}
\begin{lstlisting}[language=C, caption=Estados posibles de la alarma]
typedef enum { 
    ALARM_IDLE,      // Alarma inactiva
    ALARM_ACTIVE,    // Alarma activada/sonando
    ALARM_COOLDOWN   // Período de enfriamiento (10 segundos)
} alarm_state_t;
\end{lstlisting}
\subsection{Variables Globales}

\subsubsection{Sensores}
\begin{lstlisting}[language=C, caption=Instancias globales de sensores]
sensor_t sensor_ldr = {0.0f, 0.0f, 100.0f, 50.0f, 0, 0, 0, 0};
// valor=0, min=0%, max=100%, umbral=50%

sensor_t sensor_ntc = {0.0f, 25.0f, 30.0f, 27.5f, 0, 0, 0, 0};
// valor=0, min=25 deg C, max=30 deg C, umbral=27.5 deg C
\end{lstlisting}

\subsubsection{Control de Sensores}
\begin{lstlisting}[language=C, caption=Variables de control]
uint8_t selected_sensor = 0;        // Sensor activo: 0=LDR, 1=NTC
uint32_t last_pot_value = 0xFFFF;   // Valor anterior del potenciómetro
alarm_state_t alarm_state = ALARM_IDLE;     // Estado de la alarma
uint32_t alarm_cooldown_start = 0;  // Timestamp del inicio del cooldown
uint8_t btn_izq_last = 1;           // Estado anterior botón izquierdo
uint8_t btn_der_last = 1;           // Estado anterior botón derecho
\end{lstlisting}

\subsubsection{LEDs}
\begin{lstlisting}[language=C, caption=Mapeo de pines LED]
GPIO_TypeDef* LED_PORT[8] = {LED8_GPIO_Port, LED7_GPIO_Port, 
                             LED6_GPIO_Port, LED5_GPIO_Port, 
                             LED4_GPIO_Port, LED3_GPIO_Port, 
                             LED2_GPIO_Port, LED1_GPIO_Port};
uint16_t LED_PIN[8] = {LED8_Pin, LED7_Pin, LED6_Pin, LED5_Pin, 
                       LED4_Pin, LED3_Pin, LED2_Pin, LED1_Pin};
\end{lstlisting}

\subsection{Funciones Privadas}

\subsubsection{ADC\_ReadChannel()}
\begin{lstlisting}[language=C, caption=Lectura de canal ADC]
static uint32_t ADC_ReadChannel(uint32_t channel)
\end{lstlisting}

\textbf{Descripción:} Lee el valor digital de un canal ADC específico.

\textbf{Parámetros:}
\begin{itemize}
    \item \texttt{channel}: identificador del canal ADC a leer
\end{itemize}

\textbf{Retorna:} Valor digital de 12 bits (0-4095)

\textbf{Detalles:}
\begin{itemize}
    \item Configura el canal
    \item Inicia conversión ADC
    \item Espera a que la conversión termine
    \item Obtiene el valor y detiene el ADC
\end{itemize}

\subsubsection{Update\_Sensors()}
\begin{lstlisting}[language=C, caption=Actualización de valores de sensores]
static void Update_Sensors(void)
\end{lstlisting}

\textbf{Descripción:} Actualiza los valores de los sensores LDR, NTC y el potenciómetro.

\textbf{Funcionalidad:}
\begin{itemize}
    \item \textbf{LDR:} Convierte lectura ADC a porcentaje (0-100\%), invertida (mayor ADC = menor luminosidad)
    \item \textbf{NTC:} Calcula temperatura en $^\circ$C usando ecuación de Steinhart-Hart
    \item \textbf{Potenciómetro:} Ajusta el nivel de alarma del sensor actualmente seleccionado con histéresis de 30 unidades para evitar ruido
\end{itemize}

\subsubsection{Update\_Display()}
\begin{lstlisting}[language=C, caption=Actualización de visualización en LEDs]
static void Update_Display(void)
\end{lstlisting}

\textbf{Descripción:} Controla el encendido/apagado de LEDs según el valor del sensor activo.

\textbf{Funcionalidad:}
\begin{itemize}
    \item Calcula cuántos LEDs deben estar encendidos basándose en el valor relativo del sensor
    \item Identifica el LED que corresponde al umbral de alarma
    \item Genera parpadeo a 10 Hz en el LED del umbral
    \item Actualiza el estado de todos los LEDs
\end{itemize}

\subsubsection{Update\_Alarm()}
\begin{lstlisting}[language=C, caption=Actualización del estado de alarma]
static void Update_Alarm(void)
\end{lstlisting}

\textbf{Descripción:} Gestiona la máquina de estados de la alarma.

\textbf{Estados:}
\begin{itemize}
    \item \textbf{ALARM\_IDLE:} Monitorea si algún sensor supera su umbral
    \item \textbf{ALARM\_ACTIVE:} Alarma sonando, se mantiene hasta pulsar botón derecho
    \item \textbf{ALARM\_COOLDOWN:} Período de 10 segundos antes de poder activar nuevamente
\end{itemize}

\subsubsection{Process\_Buttons()}
\begin{lstlisting}[language=C, caption=Procesamiento de entrada de botones]
static void Process_Buttons(void)
\end{lstlisting}

\textbf{Descripción:} Detecta pulsaciones de botones y ejecuta acciones correspondientes.

\textbf{Acciones:}
\begin{itemize}
    \item \textbf{Botón Izquierdo:} Cambia el sensor activo entre LDR (0) y NTC (1)
    \item \textbf{Botón Derecho:} Desactiva la alarma y entra en período de cooldown si está activa
\end{itemize}

\textbf{Implementación:} Detecta flancos descendentes.

\subsection{Funciones Públicas}

\subsubsection{Monitor\_Init()}
\begin{lstlisting}[language=C, caption=Inicialización del módulo]
void Monitor_Init(void)
\end{lstlisting}

\textbf{Descripción:} Inicializa el módulo de monitoreo.

\textbf{Funcionalidad:}
\begin{itemize}
    \item Apaga todos los LEDs (GPIO\_PIN\_RESET)
    \item Apaga el buzzer
\end{itemize}

\textbf{Debe llamarse:} Una vez en la función \texttt{main()} durante la inicialización del sistema.

\subsubsection{Monitor\_Loop()}
\begin{lstlisting}[language=C, caption=Bucle principal de monitoreo]
void Monitor_Loop(void)
\end{lstlisting}

\textbf{Descripción:} Función principal que debe llamarse periódicamente desde el bucle principal.

\textbf{Ciclo de ejecución:} Se ejecuta cada 20 ms (50 Hz)

\textbf{Secuencia de operaciones:}
\begin{enumerate}
    \item \texttt{Process\_Buttons()}: Procesa entrada de botones
    \item \texttt{Update\_Sensors()}: Actualiza valores de sensores
    \item \texttt{Update\_Alarm()}: Gestiona estado de alarma
    \item \texttt{Update\_Display()}: Actualiza visualización en LEDs
\end{enumerate}

\textbf{Temporizador:} Utiliza \texttt{HAL\_GetTick()} para asegurar una ejecución consistente cada 20 ms.

\section{Flujo de Operación}

El sistema opera de la siguiente manera:

\begin{enumerate}
    \item \textbf{Lectura de Sensores:} El ADC lee continuamente los tres canales (LDR, NTC, Potenciómetro)
    \item \textbf{Procesamiento:} Los valores se convierten a unidades significativas (porcentaje, $^\circ$C)
    \item \textbf{Control de Alarma:} Se comparan con el umbral ajustable mediante el potenciómetro
    \item \textbf{Visualización:} Los LEDs muestran el nivel relativo, con parpadeo en el umbral
    \item \textbf{Interacción:} Los botones permiten cambiar de sensor o silenciar la alarma
\end{enumerate}

\end{document}
